\section{总结与后续工作}
当前结果支持：在本项目的文本输入与串行主机预处理基线下，利用OpenMP优化解析、编号映射和排序去重，可以在保持清洗后CSR及精确计数不变的前提下改善完整任务时间。额外内存、并行准备与串行残余仍需量化，且本研究没有完成线程绑定、NUMA和专用硬件条件下的实验。

下一阶段应先归档原始样本、输入哈希与环境元数据，再补充稳健的波动统计和峰值主存。若继续实现warp内局部调度，应让计数基线与调度优化版使用相同的32线程预处理，单独评估调度的收益和成本，并继续报告完整总时间。尚未实现的调度、逐顶点计数、外部论文复现和其他合成图实验不能作为已完成贡献写入结论。相关工作与正式参考文献仍需补充。
